\section{Objectives of the Project}
\label{sec:objectives}

\begin{itemize}
    \item \textbf{Digitize estate-to-factory operations.} \\The system transforms manual tea collection, delivery, and payment workflows into a centralized digital platform, accessible via a mobile application for field operations and a web-based admin portal for administrative and management tasks.
    \item \textbf{Reduce manual paperwork and errors.} \\Collection records, delivery confirmations, and payment deductions are stored digitally to minimize duplication and human error.
    \item \textbf{Improve verification and accountability.} \\Photographic evidence , along with structured complaint handling workflows, strengthen traceability and reduce disputes over delivery weights and fertilizer quality.
    \item \textbf{Enable fertilizer and beneficiary item workflow management.} \\The system supports fertilizer and other beneficiary item (e.g., rice, food) requests, approvals, borrowing, transport tracking, and cost deductions.
    \item \textbf{Support offline field operations.} \\The mobile application allows tea collection agents and field staff to record data offline using local storage, with automatic synchronization to the central database once connectivity is restored.
    \item \textbf{Provide transparent reporting and analytics.} \\Dashboards, charts, and graphical reports help stakeholders monitor daily and monthly collections, employee performance, stock status, and income generation.
    \item \textbf{Ensure secure and role-based access.} \\JWT-based authentication and role-based access control ensure that estate owners, managers, agents, and factory staff can only access data and functionality relevant to their role.
    \item \textbf{Ensure reliable performance under load.} \\Caching and load testing practices are incorporated to keep the system responsive as the number of users, estates, and factories grows.
    \item \textbf{Build a scalable and maintainable platform.} \\A modular, containerized architecture with automated deployment supports future expansion to other agricultural collection workflows and additional users.
    \item \textbf{Ensure software quality through structured testing.} \\Unit, end-to-end, and load testing practices are integrated into the development process to validate system reliability before deployment.
\end{itemize}